\documentclass[12pt]{article}
\usepackage[utf8]{inputenc}
\usepackage[T1]{fontenc}
\usepackage[margin=1in]{geometry}
\usepackage{setspace}
\usepackage{xcolor}
\usepackage[colorlinks=true, urlcolor=blue!60!black, linkcolor=blue!60!black]{hyperref}
\setlength{\parindent}{2em}
\setlength{\parskip}{1.25em}
\renewcommand{\baselinestretch}{1.0}
\begin{document}
\begin{center}
{\LARGE \textit{Yick Wo v. Hopkins}, 118 U.S. 356 (1886)\par}
\vspace{0.8em}
{\large Case Brief --- Four-Page Edition\par}
\end{center}
\vspace{1em}

\section{Memory Jogger}

San Francisco denied laundry permits to every Chinese applicant while granting nearly all non-Chinese applicants; a unanimous Court held that a facially neutral law administered ``with an evil eye and an unequal hand'' violates equal protection, and that the Fourteenth Amendment protects noncitizens as ``persons'' (\href{https://supreme.justia.com/cases/federal/us/118/356/}{118 U.S. 356}).

\section{Detailed Case Facts}

Amid anti-Chinese agitation in 1880, San Francisco's Board of Supervisors made it unlawful to operate a laundry without the Board's consent unless it was located in a brick or stone building, punishable by up to a \$1,000 fine and six months in jail. The ordinances set no fire-safety or sanitary criteria for consent and did not mention race or nationality.

Yick Wo, a Chinese subject who came to California in 1861, had run a laundry in the same wooden building for twenty-two years. Fire wardens and the Health Officer certified his premises safe, and he held a city license, yet the Board denied his application without reason. Of about 320 laundries, some 240 were Chinese-owned, and about 310 were in wooden buildings, like nine-tenths of the city's houses. All of the roughly 200 Chinese petitions were denied and more than 150 Chinese laundrymen were arrested, while every one of the roughly eighty non-Chinese petitions but one was granted.

\section{Procedural History}

Yick Wo was convicted, fined \$10, and jailed by Sheriff Peter Hopkins when he did not pay. The California Supreme Court denied habeas, treating the ordinances as fire-safety regulations under \textit{Barbier v. Connolly} and \textit{Soon Hing v. Crowley}. In the companion \textit{Wo Lee} case, federal Circuit Judge Lorenzo Sawyer found the ordinance plainly aimed at destroying Chinese laundries but, out of comity, remanded the petitioner so the Supreme Court could resolve the conflict. The Supreme Court decided both cases on May 10, 1886.

\section{Judicial Votes}

Reversed, \textbf{9--0}, with directions to discharge the petitioners.
\begin{itemize}
\item \textbf{Opinion of the Court:} Justice Matthews, joined by Chief Justice Waite and Justices Miller, Field, Bradley, Harlan, Woods, Gray, and Blatchford.
\item \textbf{Concurrences and dissents:} None.
\end{itemize}

\section{Holding}

(1) The Due Process and Equal Protection Clauses apply to ``all persons within the territorial jurisdiction, without regard to any differences of race, of color, or of nationality,'' including resident noncitizens. (2) A law ``fair on its face and impartial in appearance'' violates equal protection if officials administer it to make unjust discriminations between similarly situated persons by race or nationality. (3) An ordinance vesting officials with naked, arbitrary power, unbounded by objective standards, to forbid a lawful occupation is unconstitutional.

\section{Analysis of Opinions}

\textbf{Independent construction.} On federal review, the Court construes ordinances itself. These conferred not discretion based on each case's circumstances but ``a naked and arbitrary power to give or withhold consent not only as to places, but as to persons,'' exercisable ``without reason and without responsibility.''

\textbf{Distinguishing \textit{Barbier} and \textit{Soon Hing}.} Those cases upheld a uniform laundry curfew applied alike to all. These ordinances set no standard operators could meet and made them ``tenants at will, under the supervisors, of their means of living,'' unlike liquor licensing, which requires a quasi-judicial determination of fitness.

\textbf{Noncitizens as persons.} The 1880 treaty with China promised most-favored-nation protection, and the Fourteenth Amendment shifts from ``citizens'' to ``any person'' in its Due Process and Equal Protection Clauses, which are ``a pledge of the protection of equal laws.''

\textbf{Arbitrary power.} In America ``the law is the definition and limitation of power''; holding one's livelihood at another's will is ``the essence of slavery itself.'' Analogizing to voting as a right ``preservative of all rights,'' Matthews signaled that unbridled discretion is suspect.

\textbf{Discriminatory administration.} The admitted record showed the ordinances ``in actual operation'': safe Chinese operators were refused while comparable non-Chinese operators were licensed, so hostility to the petitioners' race and nationality was the only explanation.

\textbf{Later developments.} \textit{Washington v. Davis} required discriminatory purpose and recast \textit{Yick Wo} as a case of disparity so stark that intent was inescapable; \textit{Arlington Heights} called such cases ``rare''; and \textit{United States v. Armstrong} demands evidence of effect and intent even for selective-prosecution discovery.

\section{Examples: Future Applications}

\begin{enumerate}
\item \textbf{Same-Side (Invalid): Standardless Fire Permits.} A board with no criteria denies 142 of 145 applications from Central American commissary operators, most with clean inspections, while granting 38 of 40 others; noncitizens are protected, and the disparity shows an ``evil eye and an unequal hand.''
\item \textbf{Same-Side (Invalid): Selective Vending Citations.} Police issue 205 of 212 sidewalk-vending citations to Somali vendors, while other vendors documented at the same stops go uncited and no explanation is offered, supporting an inference of purpose.
\item \textbf{Opposite-Side (Valid): Uniform Safety Code.} A costly fire-suppression mandate applied mechanically to all commissaries, adopted on actuarial advice after fatal fires, closes more immigrant businesses but tracks \textit{Barbier}; impact alone is not purpose under \textit{Davis} and \textit{Feeney}.
\item \textbf{Opposite-Side (Fails): No Comparators.} A Chinese defendant prosecuted for gambling who shows only general hostility, with no unprosecuted non-Chinese comparators, loses, as in \textit{Ah Sin v. Wittman}.
\item \textbf{Fence-Sitter: Moderate Disparity.} A standardless food-truck board denies 45\% of noncitizen and 20\% of citizen applications, with inconsistent neutral reasons and one hostile remark; whether a court infers purpose is uncertain.
\end{enumerate}

\section{Critique}

\textbf{Two theories blended.} Matthews combined a structural objection to ``naked and arbitrary power'' over a lawful occupation with an equal protection objection to discriminatory enforcement; later cases such as \textit{Kotch} largely abandoned the first for economic regulation while canonizing the second.

\textbf{An evidentiary unicorn and an irony.} Because San Francisco stipulated to its lopsided record, \textit{Yick Wo} is an outlier, and using it as the benchmark under \textit{Davis}, \textit{Arlington Heights}, and \textit{Armstrong} makes subtler discrimination hard to prove. Decided four years after the Chinese Exclusion Act, it protected resident noncitizens' livelihoods even as the Court deferred to federal exclusion powers.

\section{Key Quotations}

\begin{quote}
``Though the law itself be fair on its face and impartial in appearance, yet, if it is applied and administered by public authority with an evil eye and an unequal hand, so as practically to make unjust and illegal discriminations between persons in similar circumstances, material to their rights, the denial of equal justice is still within the prohibition of the Constitution.''
\end{quote}
\noindent --- Justice Matthews, Opinion of the Court, 118 U.S. at 373--374.

\begin{quote}
``The Fourteenth Amendment to the Constitution is not confined to the protection of citizens. It says: `Nor shall any State deprive any person of life, liberty, or property without due process of law; nor deny to any person within its jurisdiction the equal protection of the laws.' These provisions are universal in their application to all persons within the territorial jurisdiction, without regard to any differences of race, of color, or of nationality, and the equal protection of the laws is a pledge of the protection of equal laws.''
\end{quote}
\noindent --- Justice Matthews, Opinion of the Court, 118 U.S. at 369.

\begin{quote}
``Though not regarded strictly as a natural right, but as a privilege merely conceded by society according to its will under certain conditions, nevertheless it is regarded as a fundamental political right, because preservative of all rights.''
\end{quote}
\noindent --- Justice Matthews, Opinion of the Court, 118 U.S. at 370.

\end{document}
